\documentclass[a4paper,fontset=windows]{ctexart}

\usepackage{anyfontsize}
\usepackage{amsmath}
\allowdisplaybreaks
\usepackage{graphicx,subfigure}
\usepackage{multirow}
\usepackage{longtable}
\usepackage{booktabs}
\usepackage{bm}
\usepackage{geometry}
\geometry{top=3cm,bottom=3cm,left=2.25cm,right=2.25cm}
\usepackage[shortlabels]{enumitem}
\usepackage{caption}
\captionsetup{font=Large}
\usepackage{indentfirst}
\setlength{\parindent}{0em}

\begin{document}\Large

    \title{\huge\textbf{实验十八\ \ 弗兰克——赫兹实验}}
    \author{\LARGE 1900011604 张植竣}
    \date{\LARGE 2020年12月10日}

    \maketitle

    \section*{\zihao{-2}1\ \  实验数据与数据处理}
    \bigskip

    \subsection*{\zihao{3}1.1\ \ 汞管的弗兰克——赫兹实验}

    \textbf{仪器号：4，$\boldsymbol{U_{\mathbf{Kg_{1}}} = 1.00\,\mathbf{V}}$，$\boldsymbol{U_{\mathbf{g_{2}p}} = 2.03\,\mathbf{V}}$，$\boldsymbol{T = 175\text{\textbf{\textcelsius}}}$，弱电流放大器的反馈电阻$\boldsymbol{R_{\mathrm{f}}}$约为$\boldsymbol{1.0\,\mathrm{M}\Omega}$。}

    测量得到的数据如下：
    \begin{center}\Large
        \begin{longtable}{ cccccc }
            \toprule[0.1em]
            $U_{\mathrm{Kg_2}}/\mathrm{V}$ & $U_{\mathrm{out}}/\mathrm{mV}$ & $U_{\mathrm{Kg_2}}/\mathrm{V}$ & $U_{\mathrm{out}}/\mathrm{mV}$ & $U_{\mathrm{Kg_2}}/\mathrm{V}$ & $U_{\mathrm{out}}/\mathrm{mV}$ \\
            \midrule[0.06em]
            0.0  &  5.1  & 10.6 & 39.5  & 21.1 & 35.5  \\
            0.2  &  5.3  & 10.8 & 24.9  & 21.4 & 25.2  \\
            0.5  &  5.4  & 11.1 & 15.6  & 21.6 & 22.3  \\
            0.7  &  5.4  & 11.2 & 13.7  & 21.7 & 20.6  \\
            0.9  &  5.5  & 11.6 & 10.4  & 21.8 & 20.0  \\
            1.2  &  5.5  & 11.9 & 10.0  & 21.9 & 21.0  \\
            1.5  &  5.5  & 12.0 & 10.0  & 22.0 & 22.5  \\
            1.7  &  5.7  & 12.1 & 10.8  & 22.2 & 25.2  \\
            2.0  &  5.7  & 12.2 & 11.7  & 22.4 & 31.4  \\
            2.2  &  5.9  & 12.3 & 13.0  & 22.6 & 40.5  \\
            2.5  &  5.9  & 12.5 & 18.0  & 22.8 & 48.6  \\
            2.8  &  6.8  & 12.8 & 23.2  & 23.1 & 61.5  \\
            3.0  &  7.3  & 13.0 & 29.9  & 23.3 & 81.5  \\
            3.2  &  8.2  & 13.2 & 41.2  & 23.5 & 98.5  \\
            3.4  &  9.2  & 13.5 & 58.7  & 23.8 & 122.0 \\
            3.7  & 12.0  & 13.7 & 73.8  & 24.1 & 165.0 \\
            3.9  & 15.3  & 14.0 & 101.8 & 24.4 & 190.0 \\
            4.1  & 18.6  & 14.3 & 117.0 & 24.5 & 195.0 \\
            4.4  & 25.3  & 14.4 & 123.8 & 24.6 & 198.5 \\
            4.7  & 33.3  & 14.5 & 129.0 & 24.7 & 204.0 \\
            4.8  & 38.0  & 14.6 & 136.4 & 24.8 & 202.0 \\
            4.9  & 41.0  & 14.7 & 142.0 & 24.9 & 193.0 \\
            5.0  & 43.7  & 14.8 & 143.8 & 25.0 & 185.6 \\
            5.1  & 45.4  & 14.9 & 135.7 & 25.2 & 169.3 \\
            5.2  & 46.8  & 15.0 & 128.9 & 25.5 & 122.0 \\
            5.3  & 45.4  & 15.3 & 96.4  & 25.7 & 113.0 \\
            5.4  & 42.3  & 15.5 & 64.8  & 25.9 & 93.0  \\
            5.6  & 33.3  & 15.8 & 37.3  & 26.1 & 67.5  \\
            5.8  & 21.9  & 16.0 & 27.3  & 26.4 & 48.7  \\
            6.1  & 14.7  & 16.3 & 17.5  & 26.7 & 40.0  \\
            6.3  & 11.2  & 16.4 & 15.3  & 26.8 & 37.8  \\
            6.6  &  9.8  & 16.5 & 14.3  & 26.9 & 36.5  \\
            6.8  &  9.0  & 16.7 & 13.0  & 27.0 & 36.0  \\
            7.0  &  8.8  & 16.8 & 12.8  & 27.1 & 37.3  \\
            7.2  &  8.5  & 16.9 & 13.5  & 27.2 & 38.8  \\
            7.3  &  8.9  & 17.0 & 14.8  & 27.4 & 41.7  \\
            7.4  &  9.7  & 17.3 & 19.0  & 27.6 & 50.2  \\
            7.5  & 10.3  & 17.5 & 23.5  & 27.9 & 63.0  \\
            7.6  & 11.0  & 17.7 & 31.2  & 28.1 & 79.0  \\
            7.9  & 16.0  & 18.0 & 46.3  & 28.4 & 110.0 \\
            8.1  & 20.0  & 18.3 & 62.9  & 28.6 & 131.5 \\
            8.3  & 26.3  & 18.6 & 91.3  & 29.0 & 183.6 \\
            8.5  & 32.6  & 18.8 & 107.0 & 29.4 & 206.5 \\
            8.8  & 45.5  & 19.0 & 128.0 & 29.5 & 211.2 \\
            9.0  & 53.3  & 19.3 & 156.5 & 29.6 & 213.0 \\
            9.3  & 69.3  & 19.5 & 171.6 & 29.7 & 211.0 \\
            9.5  & 81.6  & 19.6 & 177.1 & 29.8 & 209.4 \\
            9.7  & 93.8  & 19.7 & 178.5 & 29.9 & 206.2 \\
            9.8  & 99.0  & 19.8 & 176.8 & 30.1 & 194.8 \\
            9.9  & 101.0 & 19.9 & 172.0 & 30.4 & 171.2 \\
            10.0 & 99.0  & 20.0 & 154.6 & 30.7 & 139.9 \\
            10.1 & 95.5  & 20.3 & 121.5 & 30.9 & 123.2 \\
            10.2 & 80.8  & 20.6 & 84.0  & 31.1 & 105.1 \\
            10.3 & 75.5  & 20.9 & 51.6  & 31.3 & 86.0  \\
            \bottomrule[0.1em]
        \end{longtable}
    \end{center}

    \textbf{其它条件不变，$\boldsymbol{U_{\mathbf{g_{2}p}}}$改为$\boldsymbol{U_{\mathbf{g_{2}p}} = 3.02\,\mathbf{V}}$。}

    测量得到的数据如下：
    \begin{center}\Large
        \begin{longtable}{ cccccc }
            \toprule[0.1em]
            $U_{\mathrm{Kg_2}}/\mathrm{V}$ & $U_{\mathrm{out}}/\mathrm{mV}$ & $U_{\mathrm{Kg_2}}/\mathrm{V}$ & $U_{\mathrm{out}}/\mathrm{mV}$ & $U_{\mathrm{Kg_2}}/\mathrm{V}$ & $U_{\mathrm{out}}/\mathrm{mV}$ \\
            \midrule[0.06em]
            13.2 & 11.0 & 16.6 & 10.4 & 20.5 & 53.5 \\
            13.4 & 15.0 & 16.7 & 10.0 & 20.7 & 44.6 \\
            13.6 & 17.6 & 16.8 & 9.2  & 21.0 & 28.8 \\
            13.9 & 26.4 & 16.9 & 9.2  & 21.3 & 19.3 \\
            14.1 & 33.3 & 17.0 & 8.5  & 21.5 & 15.0 \\
            14.4 & 50.4 & 17.1 & 8.5  & 21.7 & 13.5 \\
            14.5 & 53.4 & 17.3 & 8.5  & 21.9 & 11.9 \\
            14.6 & 58.2 & 17.5 & 9.0  & 22.0 & 10.8 \\
            14.7 & 60.6 & 17.9 & 11.2 & 22.1 & 10.7 \\
            14.8 & 65.0 & 18.2 & 15.4 & 22.2 & 10.4 \\
            14.9 & 66.0 & 18.4 & 21.4 & 22.4 & 11.0 \\
            15.0 & 65.0 & 18.7 & 32.3 & 22.5 & 11.3 \\
            15.1 & 59.0 & 19.0 & 46.1 & 22.6 & 11.6 \\
            15.2 & 55.7 & 19.3 & 68.6 & 22.7 & 12.3 \\
            15.4 & 44.7 & 19.4 & 72.4 & 22.8 & 13.9 \\
            15.6 & 32.2 & 19.5 & 77.0 & 23.0 & 16.8 \\
            15.8 & 26.8 & 19.7 & 82.0 & 23.3 & 24.4 \\
            16.0 & 19.6 & 19.8 & 83.5 & 23.5 & 32.0 \\
            16.3 & 13.3 & 19.9 & 84.6 & 23.7 & 39.2 \\
            16.4 & 11.8 & 20.1 & 76.0 & 24.0 & 52.0 \\
            16.5 & 10.6 & 20.3 & 67.0 & 24.3 & 71.4 \\
            \bottomrule[0.1em]
        \end{longtable}
    \end{center}

    \newpage
    对数据点以$\Delta U = 0.001\,\mathrm{V}$为间隔进行三次插值法拟合，作图如下：

    \begin{figure}[ht!]
        \includegraphics[width = 1\linewidth]{Figure_1.png}
        \caption{汞管的弗兰克——赫兹实验}
    \end{figure}

    计算得峰值的加速电压$U_{\mathrm{Kg_2}}^{\mathrm{peak}}$与输出电压$U_{\mathrm{out}}^{\mathrm{peak}}$为：
    \begin{table}[ht!]\large
        \begin{center}
            \begin{tabular}{ |c|c|c|c|c|c|c| }
                \hline
                $i$ & 1 & 2 & 3 & 4 & 5 & 6 \\
                \hline
                $U_{\mathrm{Kg_2}}^{\mathrm{peak}}/\mathrm{V}$ & 5.206 & 9.890 & 14.774 & 19.683 & 24.729 & 29.589 \\
                \hline
                $U_{\mathrm{out}}^{\mathrm{peak}}/\mathrm{mV}$ & 46.805 & 101.038 & 144.271 & 178.563 & 204.441 & 213.032 \\
                \hline
            \end{tabular}
        \end{center}
    \end{table}

    最小二乘法拟合得到：
    \[
    U_{\mathrm{Kg_2}}^{\mathrm{peak}} = ai + b
    \]
    \begin{table}[ht!]\Large
        \begin{center}
            \begin{tabular}{ |c|c|c| }
                \hline
                $a/\mathrm{V}$ & $b/\mathrm{V}$ &      $r$       \\ \hline
                $4.895457143$  & $0.1777333333$ & $0.9999466399$ \\ \hline
            \end{tabular}
        \end{center}
    \end{table}
    \vspace{-0.5em}

    \newpage
    误差估计：

    1. 系统误差：

    计算单个测量值$U_i^{\mathrm{peak}}$的剩余方差：
    \[
    \sigma_\varepsilon^2 = \frac{1}{6}\sum_{i=1}^{6}\varepsilon_i^2 = \frac{1}{6}\sum_{i=1}^{6}\left(U_i^{\mathrm{peak}} - ai - b\right)^2 = 0.01119041905\,\mathrm{V^2}
    \]
    导致的斜率不确定度为：
    \[
    \sigma_{a,A} = a\sqrt{\dfrac{1/r^2-1}{n-2}} = 0.02528739838\,\mathrm{V}
    \]
    仪器允差$e=0.1\,\mathrm{V}$导致的不确定度：
    \[
    \sigma_{e} = \frac{e}{\sqrt{3}} = 0.02061552813\,\mathrm{V}
    \]

    2. 随机误差：

    考虑输出电压$U_{\mathrm{out}} = U_{\mathrm{out}}^{\mathrm{peak}} - 5\,\mathrm{mV}$对应的电压值$U_{i1}$、$U_{i2}$（$U_{i1}<U_{i2}$）。

    估计随机误差为
    \[
    \sigma = \frac{1}{6}\sum_{i=1}^{6}\frac{U_{i2}-U_{i1}}{2} = 0.17058333333\,\mathrm{V}
    \]

    合成的纵坐标$y = U^{\mathrm{peak}}_{\mathrm{Kg_2}}$的不确定度为：
    \[
    \sigma_{y} = \sqrt{\sigma_\varepsilon^2 + \sigma_{e}^2 + \sigma^2} = 0.42636236172\,\mathrm{V}
    \]

    则斜率$a$的不确定度为：
    \[
    \sigma_{a} = \frac{\sigma_y}{\sqrt{\sum\limits_{i=1}^{6}\left(i-\bar{i}\right)^2}} = 0.1019200985344\,\mathrm{V}
    \]

    \textbf{则最终测得第一激发态对应电位差的结果为：}
    \[
    \boldsymbol{v = 4.90\pm0.10\,\mathrm{V}}
    \]

    \newpage
    \subsection*{\zihao{3}1.2\ \ 氩管的弗兰克——赫兹实验}

    \textbf{仪器号：14，$\boldsymbol{U_{\mathbf{Kg_{1}}} = 1.00\,\mathbf{V}}$，$\boldsymbol{U_{\mathbf{g_{2}p}} = 2.03\,\mathbf{V}}$，$\boldsymbol{U_{\mathbf{F}} = 2.03\,\mathbf{V}}$。}

    测量得到的数据如下：
    \begin{center}\Large
    \begin{longtable}{ cccccccc }
        \toprule[0.1em]
        $U_{\mathrm{Kg_2}}/\mathrm{V}$ & $I/\mathrm{nA}$ & $U_{\mathrm{Kg_2}}/\mathrm{V}$ & $I/\mathrm{nA}$ & $U_{\mathrm{Kg_2}}/\mathrm{V}$ & $I/\mathrm{nA}$ & $U_{\mathrm{Kg_2}}/\mathrm{V}$ & $I/\mathrm{nA}$ \\
        \midrule[0.06em]
        0.0  & $-0.6$ & 29.0 & 32.7 & 52.2 & 63.7 & 74.5 & 74.1  \\
        1.0  & $-0.6$ & 29.5 & 29.7 & 52.4 & 63.6 & 75.0 & 77.8  \\
        2.0  & $-0.6$ & 30.0 & 25.6 & 52.6 & 63.2 & 75.5 & 81.5  \\
        3.0  & $-0.6$ & 30.5 & 21.2 & 52.8 & 62.8 & 76.0 & 84.6  \\
        4.0  & $-0.6$ & 31.0 & 15.8 & 53.0 & 62.1 & 76.5 & 87.3  \\
        5.0  & $-0.6$ & 31.5 & 10.3 & 53.5 & 60.0 & 77.0 & 90.2  \\
        5.5  & $-0.6$ & 32.0 & 5.4  & 54.0 & 56.5 & 77.5 & 92.2  \\
        6.0  & $-0.6$ & 32.5 & 2.8  & 54.5 & 52.0 & 77.7 & 92.7  \\
        6.5  & $-0.6$ & 32.7 & 2.3  & 55.0 & 47.1 & 77.9 & 93.3  \\
        7.0  & $-0.6$ & 32.9 & 1.9  & 55.5 & 40.6 & 78.1 & 93.7  \\
        7.5  & $-0.6$ & 33.1 & 1.9  & 56.0 & 34.7 & 78.3 & 94.1  \\
        8.0  & $-0.6$ & 33.3 & 2.2  & 56.5 & 29.0 & 78.5 & 94.3  \\
        8.5  &  -0.2  & 33.5 & 3.3  & 57.0 & 25.6 & 78.7 & 94.4  \\
        9.0  &  0.8   & 33.7 & 5.0  & 57.2 & 25.4 & 78.9 & 94.4  \\
        9.5  &  3.1   & 33.9 & 7.3  & 57.4 & 25.4 & 79.1 & 94.3  \\
        10.0 &  7.3   & 34.1 & 10.8 & 57.6 & 26.2 & 79.3 & 94.1  \\
        10.5 &  12.0  & 34.3 & 13.3 & 57.8 & 27.2 & 79.5 & 93.8  \\
        11.0 &  14.2  & 34.5 & 17.0 & 58.0 & 28.4 & 79.7 & 93.4  \\
        11.5 &  16.0  & 35.0 & 23.8 & 58.5 & 31.6 & 79.9 & 93.0  \\
        12.0 &  17.5  & 35.5 & 29.3 & 59.0 & 36.5 & 80.1 & 92.3  \\
        12.5 &  18.3  & 36.0 & 35.4 & 59.5 & 41.0 & 80.3 & 91.5  \\
        13.0 &  19.0  & 36.5 & 40.0 & 60.0 & 45.9 & 80.5 & 90.7  \\
        13.5 &  19.2  & 37.0 & 43.7 & 60.5 & 50.8 & 81.0 & 88.7  \\
        14.0 &  19.7  & 37.5 & 46.5 & 61.0 & 55.6 & 81.5 & 86.4  \\
        14.5 &  20.0  & 38.0 & 48.6 & 61.2 & 57.6 & 82.0 & 84.1  \\
        15.0 &  20.5  & 38.5 & 50.4 & 61.4 & 58.8 & 82.5 & 82.5  \\
        15.5 &  20.4  & 39.0 & 50.8 & 61.6 & 60.0 & 83.0 & 81.3  \\
        16.0 &  20.4  & 39.2 & 51.0 & 61.8 & 61.6 & 83.5 & 80.8  \\
        16.5 &  20.3  & 39.4 & 51.0 & 62.0 & 63.5 & 84.0 & 81.1  \\
        17.0 &  19.7  & 39.6 & 50.9 & 62.2 & 64.2 & 84.5 & 82.0  \\
        17.5 &  18.2  & 39.8 & 50.8 & 62.4 & 65.9 & 85.0 & 83.5  \\
        18.0 &  16.4  & 40.0 & 50.5 & 62.6 & 66.9 & 85.5 & 85.2  \\
        18.5 &  14.2  & 40.5 & 49.3 & 62.8 & 68.5 & 86.0 & 87.6  \\
        19.0 &  11.6  & 41.0 & 47.1 & 63.0 & 69.7 & 86.5 & 90.0  \\
        19.5 &  9.2   & 41.5 & 43.9 & 63.2 & 70.3 & 87.0 & 92.8  \\
        20.0 &  6.3   & 42.0 & 39.2 & 63.4 & 71.3 & 87.5 & 96.2  \\
        20.5 &  4.5   & 42.5 & 32.8 & 63.6 & 72.5 & 88.0 & 99.4  \\
        20.7 &  3.7   & 43.0 & 27.3 & 63.8 & 73.4 & 88.5 & 102.5 \\
        20.9 &  3.2   & 43.5 & 20.5 & 64.0 & 74.0 & 89.0 & 105.5 \\
        21.1 &  2.9   & 44.0 & 14.5 & 64.5 & 75.6 & 89.5 & 109.3 \\
        21.3 &  2.6   & 44.5 & 10.1 & 64.7 & 76.1 & 90.0 & 111.9 \\
        21.5 &  2.5   & 44.7 & 10.0 & 64.9 & 76.2 & 90.5 & 114.9 \\
        21.7 &  2.6   & 44.9 & 10.0 & 65.1 & 76.3 & 91.0 & 117.0 \\
        21.9 &  3.0   & 45.1 & 11.7 & 65.3 & 76.3 & 91.5 & 118.9 \\
        22.1 &  3.5   & 45.3 & 13.8 & 65.5 & 76.2 & 92.0 & 120.4 \\
        22.3 &  4.2   & 45.5 & 17.1 & 65.7 & 75.8 & 92.5 & 121.4 \\
        22.5 &  5.9   & 46.0 & 22.6 & 65.9 & 75.5 & 92.7 & 121.8 \\
        23.0 &  10.3  & 46.5 & 27.7 & 66.1 & 74.9 & 92.9 & 121.9 \\
        23.5 &  14.9  & 47.0 & 33.6 & 66.3 & 74.4 & 93.1 & 122.2 \\
        24.0 &  21.0  & 47.5 & 38.9 & 66.5 & 73.5 & 93.3 & 122.3 \\
        24.5 &  25.1  & 48.0 & 44.0 & 67.0 & 70.6 & 93.5 & 122.4 \\
        25.0 &  29.5  & 48.5 & 49.0 & 67.5 & 67.1 & 93.7 & 122.4 \\
        25.5 &  31.9  & 49.0 & 52.0 & 68.0 & 62.6 & 93.9 & 122.3 \\
        25.7 &  32.2  & 49.2 & 53.4 & 68.5 & 58.0 & 94.1 & 122.1 \\
        25.9 &  33.1  & 49.4 & 55.1 & 69.0 & 54.3 & 94.3 & 122.0 \\
        26.1 &  33.8  & 49.6 & 56.0 & 69.5 & 51.2 & 94.5 & 121.8 \\
        26.3 &  34.4  & 49.8 & 57.4 & 70.0 & 49.7 & 94.7 & 121.5 \\
        26.5 &  34.9  & 50.0 & 58.3 & 70.5 & 49.4 & 94.9 & 121.3 \\
        26.7 &  35.1  & 50.2 & 59.0 & 70.7 & 49.8 & 95.1 & 121.0 \\
        26.9 &  35.5  & 50.4 & 60.0 & 70.9 & 50.3 & 95.3 & 120.9 \\
        27.1 &  35.6  & 50.6 & 60.9 & 71.1 & 51.2 & 95.5 & 120.8 \\
        27.3 &  35.8  & 50.8 & 61.5 & 71.3 & 52.0 & 96.0 & 120.7 \\
        27.5 &  35.9  & 51.0 & 62.4 & 71.5 & 53.2 & 96.5 & 120.5 \\
        27.7 &  35.7  & 51.2 & 62.8 & 72.0 & 55.8 & 97.0 & 120.6 \\
        27.9 &  35.6  & 51.4 & 63.0 & 72.5 & 58.8 & 97.5 & 121.2 \\
        28.1 &  35.4  & 51.6 & 63.4 & 73.0 & 62.2 & 98.0 & 122.2 \\
        28.3 &  35.1  & 51.8 & 63.6 & 73.5 & 66.3 & 98.5 & 123.5 \\
        28.5 &  34.6  & 52.0 & 63.7 & 74.0 & 70.4 & 99.0 & 125.2 \\
        \bottomrule[0.1em]
    \end{longtable}
    \end{center}

    \newpage
    对数据点以$\Delta U = 0.001\,\mathrm{V}$为间隔进行三次插值法拟合，作图如下：

    \begin{figure}[ht!]
        \includegraphics[width = 1\linewidth]{Figure_2.png}
        \caption{氩管的弗兰克——赫兹实验}
    \end{figure}

    计算得峰值的加速电压$U_{\mathrm{Kg_2}}^{\mathrm{peak}}$与输出电压$I^{\mathrm{peak}}$为：
    \begin{table}[ht!]\large
        \begin{center}
            \begin{tabular}{ |c|c|c|c|c|c|c|c| }
                \hline
                $i$ & 1 & 2 & 3 & 4 & 5 & 6 & 7 \\
                \hline
                $U_{\mathrm{Kg_2}}^{\mathrm{peak}}/\mathrm{V}$ & 15.246 & 27.713 & 39.649 & 52.555 & 65.77 & 79.515 & 94.442 \\
                \hline
                $I^{\mathrm{peak}}/\mathrm{nA}$ & 20.523 & 35.908 & 51.021 & 63.709 & 76.311 & 94.413 & 122.415 \\
                \hline
            \end{tabular}
        \end{center}
    \end{table}

    最小二乘法拟合得到：
    \[
    U_{\mathrm{Kg_2}}^{\mathrm{peak}} = ai + b
    \]

    \begin{table}[ht!]\Large
        \begin{center}
            \begin{tabular}{ |c|c|c| }
                \hline
                $a/\mathrm{V}$ & $b/\mathrm{V}$ & $r$ \\
                \hline
                $13.11832143$ & $1.082428571$ & $0.9993666047$ \\
                \hline
            \end{tabular}
        \end{center}
    \end{table}

    \newpage
    误差估计：

    1. 系统误差：

    计算单个测量值$U_i^{\mathrm{peak}}$的剩余方差：
    \[
    \sigma_\varepsilon^2 = \frac{1}{5}\sum_{i=1}^{7}\varepsilon_i^2 = \frac{1}{5}\sum_{i=1}^{7}\left(U_i^{\mathrm{peak}} - ai - b\right)^2 = 1.221974507143\,\mathrm{V^2}
    \]
    导致的斜率不确定度为：
    \[
    \sigma_{a,A} = a\sqrt{\dfrac{1/r^2-1}{n-2}} = 0.20890655012\,\mathrm{V}
    \]
    仪器允差$e=0.1\,\mathrm{V}$导致的不确定度：
    \[
    \sigma_{e} = \frac{e}{\sqrt{3}} = 0.02061552813\,\mathrm{V}
    \]

    2. 随机误差：

    考虑输出电流$I = I^{\mathrm{peak}} - 0.5\,\mathrm{nA}$对应的电压值$U_{i1}$、$U_{i2}$（$U_{i1}<U_{i2}$）。

    估计随机误差为
    \[
    \sigma = \frac{1}{7}\sum_{i=1}^{7}\frac{U_{i2}-U_{i1}}{2} = 0.75650000000\,\mathrm{V}
    \]
    合成的$U_{\mathrm{Kg_2}}^{\mathrm{peak}}$的不确定度为：
    \[
    \sigma'_{U_{\mathrm{Kg_2}}^{\mathrm{peak}}} = \sqrt{\sigma_\varepsilon^2 + \sigma_{e}^2 + \sigma^2} = 1.40776696952\,\mathrm{V}
    \]

    则斜率$a$的不确定度为：
    \[
    \sigma_{a} = \frac{\sigma_y}{\sqrt{\sum\limits_{i=1}^{7}\left(i-\bar{i}\right)^2}} = 0.26604295038\,\mathrm{V}
    \]

    \textbf{则最终测得第一激发态对应电位差的结果为：}
    \[
    \boldsymbol{v = 13.12\pm0.27\,\mathrm{V}}
    \]

    \newpage
    \section*{\zihao{-2}2\ \  思考题}

    1. 可以从图1中看出：反向电压$U_{\mathrm{g_{2}p}}$增大时，图像峰值、谷值都减小，且峰、谷值对应的加速电压都稍增大一些。以峰值为例：

    \begin{table}[ht!]\large
        \begin{center}
            \begin{tabular}{ |c|c|c|c| }
                \hline
                \multirow{2}{*}{$U_{\mathrm{g_{2}p}} = 2.03\,\mathrm{V}$} & $U_{\mathrm{Kg_2}}^{\mathrm{peak}}/\mathrm{V}$ & 14.774 & 19.683 \\
                \cline{2-4}
                 & $U_{\mathrm{out}}^{\mathrm{peak}}/\mathrm{mV}$ & 144.271 & 178.563 \\
                \hline
                \multirow{2}{*}{$U_{\mathrm{g_{2}p}} = 3.02\,\mathrm{V}$} & $U_{\mathrm{Kg_2}}^{\mathrm{peak}}/\mathrm{V}$ & 14.934 & 19.886 \\
                \cline{2-4}
                 & $U_{\mathrm{out}}^{\mathrm{peak}}/\mathrm{mV}$ & 66.070 & 84.657 \\
                \hline
            \end{tabular}
        \end{center}
    \end{table}\vspace{-1.2em}

    推测是因为反向电压增大后，阻碍动能不足的电子更多，到达板极的电子减少，则电流减小，使得图像中电流的峰值、谷值减小。且需要更大的加速电压加速电子以克服反向电压的作用，导致峰、谷值对应的加速电压稍增大一些。

    \section*{\zihao{-2}3\ \  分析与讨论}

    1. 读数时$U_\mathrm{out}$和$I$不稳定，尤以汞管的$U_\mathrm{out}$为甚。时常有$U_\mathrm{out}$示数突然增大或突然减小的情况出现，示数较为稳定时$U_\mathrm{out}$会在一定范围内波动，记录的是估计的平均水平。$U_\mathrm{out}$和$I$的示数会随着时间的推移而减小，故尽量调节时让$U_{\mathrm{Kg_2}}$快速单调增大而不重复测量，否则重测的值会与之前测量的趋势不符。

    2. 在汞管的弗兰克——赫兹实验中，外接数字电压表（以$\mathrm{mV}$为单位）的示数与仪器的$I$档位示数（以$\mathrm{nA}$为单位）基本相同，可以得到弱电流放大器的反馈电阻$R_{\mathrm{f}}$约为$1.0\,\mathrm{M}\Omega$。但接数字电压表可以获得多一位有效数字。

    3.汞管的伏安特性曲线比氩管的尖锐、陡峭，而氩管输出电流的峰——谷值差距随加速电压的增大而显著减小，在$U_{\mathrm{Kg_2}} > 90\,\mathrm{V}$后峰——谷值差距几乎为0，但是实验中并没有测出比第一激发态更高的能级。

\end{document}